This thesis develops first-order optimization and online learning methods
with explicit performance guarantees for nonsmooth convex optimization,
storage control, and transmission-rate adaptation. The optimization methods
avoid costly domain projections. The storage methods account for the effect
of current decisions on future resource availability. Decision
parameterizations and auxiliary variables expose convex structure and make
the algorithmic updates tractable. For rate adaptation from binary acknowledgments,
confidence bounds guide transmission-rate selection and yield bounds on
time-average expected queue length.

For nonsmooth convex optimization, methods combining linear optimization over
the original domain with subgradient queries and projections on an auxiliary
set achieve an $O(T^{-1/2})$ upper bound on the objective gap using at most $T$
queries to each oracle. For problems with additional convex inequality
constraints, the thesis applies the algorithm published jointly with Michael
J. Neely in 2024 to the positive parts of the constraints, with a parameter
choice and analysis for exact deterministic oracles.
The constraint-violation norm decreases at the same rate when an optimal Lagrange
multiplier for the functional constraints exists. The algorithm does not
require the multiplier or a bound on its norm as input.

For energy-harvesting allocation across $n\ge2$ subsystems, a controller
updates total usage and allocation proportions while respecting energy
availability. For a planned horizon $T$ and sufficient battery capacity of
order $\sqrt T$, independent of $n$, it achieves
$O(\sqrt{T(1+\log n)})$ expected cumulative
regret against the best fixed nonnegative allocation in hindsight whose
total usage does not exceed the mean energy arrival. Arrivals and total
per-slot usage lie in $[0,1]$. Costs are convex and revealed after allocation,
with a fixed componentwise bound on the selected cost subgradients. The
arrivals are i.i.d., and their next-slot conditional mean, given the full
past history, equals their common mean.

For scalar storage control with a known retention coefficient $\alpha\in(0,1)$,
sufficient capacity to avoid saturation, bounded adversarial arrivals, and
convex costs, simplex disturbance-action control (SDAC) policies specify
allocations of past arrivals over time.
An entropic SDAC controller learns these allocation
fractions while respecting resource availability. It achieves
$O(\sqrt{T\log(H+1)})$ regret against fixed SDAC policies with
memory $H$ and $O(\sqrt{T\log\log T})$ regret against fixed infinite-memory
SDAC policies. Each comparator is evaluated on its own feasible trajectory
under the same arrivals and costs. These bounds assume full cost feedback,
with $L_x,L_u>0$ bounding the state and action components of the selected
cost subgradients.
For each memory length $H\ge1$ and all sufficiently large $T$, every causal
feasible controller has worst-case
expected regret at least
$\frac{\alpha}{8}\bigl(L_x/(1-\alpha)+L_u\bigr)\sqrt T$
against the fixed memory-$H$ SDAC class. For fixed $H,\alpha,L_x,L_u$,
the $\sqrt T$ dependence on the horizon $T$ is optimal.

For an initially empty communication queue with independent i.i.d. arrival and
channel-capacity sequences in $[0,1]$, a transmitter chooses rates from a
continuous interval and learns their mean offered service from binary
acknowledgments. Given a known positive lower bound $\varepsilon$ on the
difference between the best achievable mean offered service and the mean
arrival rate, UCB1 on a finite rate grid achieves time-average expected queue
length $O((1+\log(1/\varepsilon))/\varepsilon^2)$, uniformly over all horizons.
As $\varepsilon\downarrow0$, the converse result from the original joint
paper on rate adaptation gives a worst-case lower bound of $\Omega(\varepsilon^{-2})$
for every causal policy in some admissible environment and at some horizon.
For this UCB1
controller, a matching $\Omega(\varepsilon^{-2}\log(1/\varepsilon))$ lower
bound holds at some horizon even with deterministic arrivals and capacities.
